\documentclass[11pt,a4paper]{article}
\usepackage{cmap}
\usepackage[utf8]{inputenc}
\usepackage[T1]{fontenc}
\usepackage{lmodern}
\usepackage{microtype}
\usepackage[margin=1.0in]{geometry}
\usepackage{amsmath,amssymb,amsthm,mathtools}
\usepackage{booktabs,array,tabularx}
\usepackage[svgnames]{xcolor}
\usepackage{listings}
\usepackage[obeyspaces,spaces]{url}
\usepackage[colorlinks=true,linkcolor=DarkBlue,citecolor=DarkBlue,urlcolor=DarkBlue]{hyperref}

\urlstyle{tt}
\def\UrlFont{\ttfamily\upshape}
\renewcommand{\arraystretch}{1.3}
\setlength{\tabcolsep}{5pt}

\sloppy
\setlength{\emergencystretch}{4em}
\hfuzz=4.0pt
\vfuzz=4.0pt

\newtheorem{theorem}{Theorem}[section]
\newtheorem{lemma}[theorem]{Lemma}
\newtheorem{proposition}[theorem]{Proposition}
\newtheorem{corollary}[theorem]{Corollary}
\theoremstyle{definition}
\newtheorem{definition}[theorem]{Definition}
\theoremstyle{remark}
\newtheorem{remark}{Remark}[section]

\lstdefinestyle{lean}{
  basicstyle=\ttfamily\footnotesize, breaklines=true, frame=single,
  framesep=4pt, xleftmargin=4pt,
  literate={≤}{{$\leq$}}1 {∀}{{$\forall$}}1 {→}{{$\to$}}1 {↔}{{$\leftrightarrow$}}1
           {ℚ}{{$\mathbb{Q}$}}1 {α}{{$\alpha$}}1 {¬}{{$\neg$}}1 {·}{{$\cdot$}}1
           {∩}{{$\cap$}}1 {∅}{{$\emptyset$}}1 {∨}{{$\vee$}}1,
}
\lstset{style=lean}

\title{\textbf{A Proved Lift That Does Not Transfer:\\
Split-Level Certificates of Non-Tree-Likeness,\\
and Three Mechanisms of Failure}}
\author{Xavier Callens\\
\small SocrateAI Scientific Research \& Formal Verification Laboratory\\
\small Formal artifact: \texttt{ANSE.GeneticaReticulataTree} (Lean 4)}
\date{Edition 1, 23 October 2026. \textbf{Not peer reviewed.}}

\begin{document}
\maketitle

\begin{abstract}
A quartet certificate refutes the four-point condition on four sequences with an
exact integer comparison a proof kernel can recheck. It names four sequences;
what a biologist reads is a \emph{split}, one group of sequences against
another. We promote the certificate to that level and report both halves of the
result honestly.

The formal half succeeds. Because complete deletion gives every quadruple its own
denominator, the Bandelt--Dress isolation index --- a \emph{minimum} over
straddling quadruples --- is not available without an incommensurable
comparison, so we define the split certificate as a \emph{logical conjunction} of
per-quadruple integer comparisons instead. We then prove that two certified
\emph{incompatible} splits force a four-point violation, with the witness
quadruple \emph{constructed} from the four side-intersections rather than
searched for (four theorems; the Lean module is 27/27 declarations with no
\texttt{sorry} and no axioms beyond \{\texttt{propext},
\texttt{Classical.choice}, \texttt{Quot.sound}\}).

The empirical half fails, and the failures are the contribution. Reaching a
working simulated demonstration (20/20 specificity, 15/15 sensitivity across
seven populations at five seeds) required three negative controls to fail first:
the criterion had no noise floor and fired on a two-count margin in two thousand
where an exact tree forces equality; no constant floor can separate, because
saturation retains exactly $n-3$ splits at a floor $64\times$ the one that erases
every other dataset; and the matched null was never matched, because a
relaxed-clock simulator's divergence at fixed mutation rate has a coefficient of
variation near $50\%$ spanning $10$--$28\times$, so matching it \emph{in
expectation} is ill-posed however carefully implemented. Rejection sampling on
each replicate's \emph{realised} divergence repairs it.

The demonstration then does not transfer. At sixteen taxa on real viral
populations the strict conjunction certifies $0$, $2$ and $4$ splits where a tree
has $13$: every population reads as a tree, so its single agreement with a
tree-like prior is vacuous. The mechanism is measured, not inferred --- conjuncts
for a balanced split number $\binom{n/2}{2}^2$, which is $36$ at $n=8$ but $784$
at $n=16$, and the closest-missing splits failed by \emph{one} conjunct out of
$91$. An exception budget calibrated against the matched null recovers power and
does not separate: all three populations, including the only tree-like control,
stop reading as trees at a smaller budget than a genuine tree does. Any threshold
between $13$ and $24$ would classify all three correctly, which on three
populations with one negative control is fitting a cutoff to the answer, so we
choose none.

We conclude that the binding constraint is not the statistic but the alignment
admissibility gate, which leaves three of ten populations usable with a single
negative control. One negative control cannot separate ``this method
false-positives'' from ``that population's prior is wrong.'' The next step for
this line is more admissible negative controls, not another statistic.
\end{abstract}

\section{Introduction and what is new}\label{sec:intro}

The four-point condition~\cite{buneman} is the tree-level null for a distance on
a finite set: a metric is a tree metric exactly when, for all $a,b,c,e$,
\begin{equation}\label{eq:fourpoint}
  d(a,b) + d(c,e) \;\le\; \max\bigl(d(a,c)+d(b,e),\; d(a,e)+d(b,c)\bigr).
\end{equation}
A \emph{quartet certificate} is four sequences, their six pairwise counts, and
the one common denominator those counts share under complete deletion within the
quadruple, such that \eqref{eq:fourpoint} fails as an integer comparison. Prior
work established such certificates, measured the cost of the alignment step that
produces the columns, and showed that a certificate's margin is exactly twice one
Bandelt--Dress isolation index, so that a certificate \emph{is} the
split-decomposition condition for a split system to fail to be a tree at the
level of four taxa~\cite{paperf}.

That identity invites an obvious promotion. A quadruple is not a biological
object; a \emph{split} is. This paper asks whether the certificate lifts, and
reports three things that are new:

\begin{enumerate}
\item \textbf{A definition forced by exactness} (\S\ref{sec:conjunction}). The
isolation index of a split is a minimum over the quadruples straddling it. Under
complete deletion within the subset each such quadruple has a different
denominator, so that minimum compares incommensurable quantities. Defining the
split certificate as a \emph{conjunction} of per-quadruple integer comparisons
avoids this entirely and needs no common metric.
\item \textbf{A lift theorem} (\S\ref{sec:lift}). Two certified incompatible
splits force a four-point violation, and the witness quadruple is constructed
from the four side-intersections rather than searched for. Four machine-checked
theorems, one of which proves a claim earlier editions of~\cite{paperf} asserted
without proof.
\item \textbf{A measured non-transfer} (\S\S\ref{sec:nontransfer}--\ref{sec:three}).
The method works on one simulated grid and has no power at the taxon count real
data has. We give the scaling law responsible, the near-miss distribution that
distinguishes ``too strict'' from ``wrong,'' and the reason three populations
cannot settle the question.
\end{enumerate}

Along the way \S\ref{sec:nullmatching} reports a methodological finding with
scope beyond this paper: a matched null calibrated on the \emph{mean} of a
nuisance parameter is not matched when the nuisance has a heavy-tailed realised
distribution, and we give the measurement that detects this.

\section{A conjunction, not a minimum}\label{sec:conjunction}

Write a quadruple's three pairing sums, over its own common columns, as $s_1 \le
s_2 \le s_3$. Each of the three splits of the quadruple carries a Bandelt--Dress
isolation index~\cite{bandeltdress}
\begin{equation}\label{eq:iso}
  \alpha(A|B) \;=\; \tfrac12 \min_{a,a' \in A,\; b,b' \in B}
    \max\{\,0,\; \mathrm{cross}_1 - \mathrm{own},\; \mathrm{cross}_2 - \mathrm{own}\,\},
\end{equation}
which on a single quadruple reduces to
$\tfrac12\max\bigl(0, \max(\mathrm{cross}_1,\mathrm{cross}_2) - \mathrm{own}\bigr)$,
and whose three values are therefore $\tfrac12(s_3-s_1)$, $\tfrac12(s_3-s_2)$ and
$0$. A split system is displayed by a tree exactly when it is pairwise
compatible~\cite{semplesteel}, so a split system containing two incompatible
splits of positive index is the classical statement of non-tree-likeness.

The minimum in~\eqref{eq:iso} is the obstruction. Complete deletion
\emph{within the subset} gives each straddling quadruple its own set of columns
at which all four sequences carry a definite base, and hence its own
denominator. Minimising a margin across quadruples with different denominators
compares quantities that are not commensurable, and no choice of common
denominator is available that does not reintroduce the columns complete deletion
was there to remove.

\begin{definition}[Split certificate]\label{def:split}
Let $\theta$ be a non-negative rational. The split $A|B$ is \emph{certified at
floor $\theta$} when, for \textbf{every} quadruple $\{a,a',b,b'\}$ with $a,a' \in
A$ and $b,b' \in B$, the split's own pairing sum lies strictly below the larger
crossing sum by more than $\theta$ times that quadruple's own denominator:
\[
  \bigl(\max(\mathrm{cross}_1,\mathrm{cross}_2) - \mathrm{own}\bigr) \cdot \theta_{\mathrm{den}}
  \;>\; \theta_{\mathrm{num}} \cdot \mathrm{common}.
\]
\end{definition}

Every conjunct is an integer comparison over its own denominator; nothing is
averaged and no common metric is required. The split-level statement inherits its
exactness from the quadruple level rather than paying for it. A split isolating a
single taxon has no straddling quadruple, so an unguarded conjunction over its
quadruples would be vacuously true; such splits carry no evidence and are
excluded.

\section{The lift theorem}\label{sec:lift}

Two splits are \emph{incompatible} when all four of $A \cap C$, $A \cap D$, $B
\cap C$, $B \cap D$ are non-empty. Choosing one taxon from each gives a quadruple
on which the two splits induce two \emph{different} pairings.

\begin{theorem}[Two certified pairings force the third strictly largest]
\label{thm:lift}
Let $p$ and $q$ be the pairing sums of two different splits of a quadruple and
$r$ the third. If $p < \max(q,r)$ and $q < \max(p,r)$, then $\max(p,q) < r$, and
consequently $r \not\le \max(p,q)$: the four-point condition fails on that
quadruple.
\end{theorem}

\begin{proof}
By $\max(p,q) < r \iff p<r \wedge q<r$ it suffices to establish both. From
$p < \max(q,r)$ either $p<q$ or $p<r$; from $q < \max(p,r)$ either $q<p$ or
$q<r$. The case $p<q$ with $q<p$ is impossible. If $p<q$ and $q<r$ then $p<r$ by
transitivity and both hold. If $p<r$ and $q<p$ then $q<r$ and both hold. If $p<r$
and $q<r$ both hold directly.
\end{proof}

\begin{corollary}
If two incompatible splits are both certified, then the four-point condition
fails on the quadruple obtained by taking one taxon from each of the four
side-intersections. The witness is \emph{constructed}, not searched for.
\end{corollary}

So a population-level claim --- two groupings both certified, which no tree can
display together --- forces a quadruple-level certificate that a kernel can
recheck. The following are machine-checked in
\texttt{ANSE.GeneticaReticulataTree}:

\begin{center}\footnotesize
\begin{tabular}{@{}ll@{}}
\toprule
declaration & content \\
\midrule
\texttt{third\_pairing\_strictly\_largest\_of\_two\_certified}
  & Theorem~\ref{thm:lift} \\
\texttt{fourPoint\_violated\_of\_two\_certified}
  & the same, as a four-point violation \\
\texttt{fourSet\_splits\_pairwise\_incompatible}
  & no two of a four-set's three splits are compatible \\
\texttt{splitsCompatible\_self}
  & a split never crosses itself \\
\bottomrule
\end{tabular}
\end{center}

The module compiles 27 of 27 declarations with no \texttt{sorry} and no axioms
beyond $\{\texttt{propext}, \texttt{Classical.choice}, \texttt{Quot.sound}\}$,
checked by compiling and then reading \texttt{\#print axioms} for every
declaration, because \texttt{sorry} and a smuggled axiom both exit zero.
\texttt{fourSet\_splits\_pairwise\_incompatible} supplies a claim that editions~5
and~9 of~\cite{paperf,textbook} asserted in prose without proof.

\begin{remark}
Theorem~\ref{thm:lift} uses no tree, no complex and no metric space --- only
three rationals. That is why the quantity survives per-subset evaluation, where a
complex built over a global distance matrix cannot restrict to a subset without
reintroducing the deleted columns.
\end{remark}

\section{The simulated demonstration, and the three controls that failed first}
\label{sec:simulated}

Seven populations at five seeds, each $2000$ columns: a clonal strict-clock tree
(D1), a clonal relaxed clock (D2), one planted recombinant under each clock
(D3, D4), reciprocal mosaics (D5), gap blocks (D6), and a saturated population
(D7). D3, D4 and D5 carry reticulation; the rest do not and are the controls. We
report the three control failures because each located something the original
question did not.

\subsection{The criterion had no noise floor}\label{sec:nofloor}

With $\theta = 0$, Definition~\ref{def:split} is $\mathrm{own} <
\max(\mathrm{cross})$, and the negative control failed $0$ of $5$: the clonal
strict-clock population certified $8$ splits where its tree has $5$, and $6$ of
the resulting pairs were incompatible. The first exhibit shows why. Its pairing
sums were $[66, 110, 112]$ over a denominator of $2000$ --- and on an \emph{exact}
tree metric the two crossing sums are \textbf{equal}. Their observed difference
is sampling noise of order $\sqrt{L}$, and the sign of that noise is what the
criterion was reading: the split with $110$ was certified on a two-count margin
in two thousand.

\begin{remark}
A pre-registered prediction also failed to be a test. It asked whether a planted
recombinant was implicated, accepting it either in a witness quadruple or
anywhere in the incompatible pair's splits --- and a split's two sides cover
\emph{every} taxon, so the second disjunct is true whenever any incompatible pair
exists. Only the witness quadruple, four taxa of nine, carries information. The
forced form recorded a pass; the repaired form is what \S\ref{sec:rejection}
reports.
\end{remark}

\subsection{No constant floor can separate}\label{sec:noconstant}

Sweeping $\theta$ over ten values, no floor separates the reticulate populations
from the controls. Tree-readings per five seeds, where a tree-reading is correct
for D1, D2, D6, D7 and a miss for D3, D4, D5:

\begin{center}\footnotesize
\begin{tabular}{@{}rccccccc@{}}
\toprule
$\theta$ & D1 & D2 & D6 & D7 & D3 & D4 & D5 \\
\midrule
$0.00100$ & 5/5 & 5/5 & 1/5 & 1/5 & 0/5 & 0/5 & 0/5 \\
$0.00400$ & 5/5 & 5/5 & 5/5 & 3/5 & 2/5 & 2/5 & 0/5 \\
$0.01600$ & 5/5 & 5/5 & 5/5 & 5/5 & 5/5 & 5/5 & 5/5 \\
\bottomrule
\end{tabular}
\end{center}

The saturated population retains exactly $n-3 = 5$ certified splits at
$\theta = 0.064$, sixty-four times the floor that erases every other dataset,
because saturation inflates margins \emph{in rate terms}. A single global rate is
the wrong functional form.

\subsection{The matched null was never matched}\label{sec:nullmatching}

Calibrating the floor per population against a clonal null matched on divergence
fixed specificity but reported a null whose divergence missed the population's by
$77\%$, and only on the two relaxed-clock datasets. Two defects compounded.

First, the tuner bisected a \emph{single realisation}. The relaxed-clock
simulator draws its per-branch rate multipliers from the same pseudorandom stream
as its mutations, and the mutation step consumes a number of draws that depends
on the rate; changing the rate therefore shifts the stream and changes every
downstream multiplier. A single realisation's divergence is consequently a
chaotic, non-monotone function of the rate, and bisecting it wanders rather than
converging. The strict-clock simulator has nothing downstream of its mutations,
which is exactly why only relaxed-clock nulls were affected.

Averaging the objective over five seeds barely helped --- worst relative error
$147.5\%$ against $143.7\%$ --- so we measured the assumption underneath. At
\emph{fixed} rate, over $200$ realisations:

\begin{center}\footnotesize
\begin{tabular}{@{}rrrrr@{}}
\toprule
rate & mean divergence & CV & min & max/min \\
\midrule
$0.001$ & $0.00658$ & $57.6\%$ & $0.00088$ & $28.4\times$ \\
$0.004$ & $0.02677$ & $49.1\%$ & $0.00521$ & $13.6\times$ \\
$0.016$ & $0.09299$ & $48.9\%$ & $0.02272$ & $10.1\times$ \\
\bottomrule
\end{tabular}
\end{center}

Each lineage's rate is its parent's times $\exp(\sigma u)$, so rate multipliers
\emph{compound} across generations and their log-variance grows with depth. The
distributions at rates $0.001$ and $0.004$ overlap heavily. This is the second
and deeper defect: the floor is a \textbf{maximum} over replicates, so when the
realised divergences of those replicates differ by up to $28\times$, the maximum
is set by whichever replicate drew a fast-rate tree. The null was inflated, hence
conservative, but not matched.

\begin{remark}[A finding with scope beyond this paper]
Matching a null on the \emph{mean} of a nuisance parameter is not matching it
when the nuisance has a heavy-tailed realised distribution, and averaging more
replicates does not repair this: averaging shrinks the standard error of the
\emph{tuning}, not the spread of the \emph{replicates} the statistic is computed
on. The diagnostic is cheap and we recommend it generally --- measure the
across-realisation coefficient of variation of the quantity being matched, and if
it is comparable to the range being matched over, expectation-matching is
ill-posed.
\end{remark}

\subsection{Rejection sampling on the realised nuisance}\label{sec:rejection}

A replicate counts only if \emph{its own} realised divergence is within $10\%$ of
the population's. With that, four pre-registered predictions hold at full count:
every non-reticulate population reads as a tree ($20/20$); every reticulate
population does not ($15/15$); the reticulate ones still certify something, so an
empty system cannot masquerade as a verdict ($15/15$); and a planted recombinant
appears in a witness quadruple ($15/15$).

\begin{center}\footnotesize
\begin{tabular}{@{}llrrcrr@{}}
\toprule
dataset & reticulate & floor & certified & $n-3$ & trees & acceptance \\
\midrule
D1 clonal, strict clock & no & $0.00170$ & $\mathbf{5.0}$ & 5 & 5/5 & $72.5\%$ \\
D2 clonal, relaxed clock & no & $0.00120$ & $4.4$ & 5 & 5/5 & $17.1\%$ \\
D6 gap blocks & no & $0.00339$ & $4.0$ & 5 & 5/5 & $52.5\%$ \\
D7 saturated & no & $\mathbf{0.04000}$ & $\mathbf{5.0}$ & 5 & 5/5 & $100.0\%$ \\
D3 one recombinant & \textbf{yes} & $0.00160$ & $9.0$ & 6 & 0/5 & $54.7\%$ \\
D4 relaxed + recombinant & \textbf{yes} & $0.00190$ & $7.4$ & 6 & 0/5 & $15.5\%$ \\
D5 reciprocal mosaics & \textbf{yes} & $0.00170$ & $10.6$ & 7 & 0/5 & $51.5\%$ \\
\bottomrule
\end{tabular}
\end{center}

D1 and D7 recover exactly $n-3 = 5$ certified splits, their true internal split
count, with no incompatible pairs; D7's floor is $23\times$ D1's, which is the
divergence dependence no constant rate could express. Acceptance rates are
reported because they matter: a population that cannot be matched receives no
floor and no verdict rather than a floor from an unmatched null. Across all
$350$ dataset--seed--floor combinations, every constructed witness really
violates the four-point condition.

\section{Non-transfer at sixteen taxa}\label{sec:nontransfer}

The simulated grid has $8$--$10$ taxa. Real viral populations, filtered by
pre-registered quality criteria and aligned to a chosen reference, have $16$. The
same pipeline that produced the quartet results of~\cite{paperf} --- the same
reference-anchored alignment, the same two references, and the same
\emph{measured} admissibility threshold above which a population is reported
blocked rather than interpreted --- admits three of ten populations.

\subsection{Making the computation possible}

Certifying every split of $16$ taxa requires about $39.7$ million straddling
quadruple evaluations. There are only $\binom{16}{4} = 1820$ distinct quadruples,
each reused about $21\,800$ times, so we precompute their counts once (about
$1.3 \times 10^7$ column comparisons) and the rest are table lookups; done
directly it is about $2.8 \times 10^{11}$. The table must identify which of the
three canonical pairing sums a split induces, which it does by asking whom the
smallest index is paired with --- index arithmetic that is easy to get subtly
wrong and invisible in aggregate output, so it is checked exhaustively against
directly computed counts for every quadruple and every $2{+}2$ split, and the
refactor is confirmed to reproduce the simulated artifact exactly.

\subsection{Zero power}

A $16$-taxon tree has $n - 3 = 13$ internal splits.

\begin{center}\footnotesize
\begin{tabular}{@{}llrl@{}}
\toprule
population & prior & certified splits & reads as \\
\midrule
Coxsackievirus A6 & reticulate & $\mathbf{0}$ / 32751 & tree \\
Hepatitis B & reticulate & $\mathbf{2}$ / 32751 & tree \\
JC polyomavirus & tree-like & $\mathbf{4}$ / 32751 & tree \\
\bottomrule
\end{tabular}
\end{center}

Every population reads as a tree. The single agreement with a tree-like prior is
therefore \emph{vacuous}: a method that always answers ``tree'' agrees with every
tree-like prior for no reason, and reporting $1/3$ here would be reporting an
artefact of having no power.

\subsection{The mechanism, measured}

Definition~\ref{def:split} requires every straddling quadruple to clear the
floor, and the number of conjuncts for a split with $k$ taxa on one side is
$\binom{k}{2}\binom{n-k}{2}$, maximised at $\binom{n/2}{2}^2$: \textbf{$36$ at
$n = 8$, but $784$ at $n = 16$}. A conjunction over hundreds of noisy comparisons
dies on a single dissent.

The near-miss distribution separates ``the method is wrong about which splits
exist'' from ``the method is too strict to certify one,'' and these demand
opposite conclusions:

\begin{center}\footnotesize
\begin{tabular}{@{}lrr@{}}
\toprule
population & closest miss & conjunction length \\
\midrule
Coxsackievirus A6 & \textbf{1} failing conjunct & 91 \\
Hepatitis B & 2 failing conjuncts & 91 \\
JC polyomavirus & 11 failing conjuncts & 234 \\
\bottomrule
\end{tabular}
\end{center}

One dissenting quadruple out of ninety-one. The rule is not wrong about the
splits; it is too strict to certify one at this conjunction length, and it
succeeded on the simulated grid only because eight taxa made the requirement
roughly twenty times weaker. We had stated the simulated scope as ``one parameter
grid, $8$--$10$ taxa'' and treated that as adequate; it named parameters without
identifying which one the method's difficulty scales in, and taxon count was that
one.

\section{An exception budget, and why three populations cannot settle it}
\label{sec:three}

Allowing at most $k$ dissenting conjuncts keeps the condition \emph{exact and
countable} --- ``at most $k$ of these integer comparisons fail'' is as checkable
as ``none does'' --- so it is a weaker claim, not a vaguer one. What it must never
be is a chosen number. We calibrate $k$ as the largest value at which
\emph{every} matched null replicate still reads as a tree: one step beyond, a
genuine clonal tree at the population's own divergence would be called
reticulate. This gives $k = 23$, with the null first breaking at $24$.

That recovers power and produces a false positive: Coxsackievirus A6 certifies
$8$ splits and reads as a non-tree (correct), Hepatitis B certifies $18$ and
reads as a non-tree (correct), and JC polyomavirus certifies $6$ and reads as a
\textbf{non-tree}, which its prior says is wrong.

A threshold invites fitting, so we replace it with a comparison carrying no
cutoff: the budget at which the population stops reading as a tree, against the
budget at which its \emph{own} matched null does.

\begin{center}\footnotesize
\begin{tabular}{@{}llrrr@{}}
\toprule
population & prior & population breaks at & null breaks at & difference \\
\midrule
Hepatitis B & reticulate & 3 & 24 & $-21$ \\
Coxsackievirus A6 & reticulate & 9 & 24 & $-15$ \\
JC polyomavirus & \textbf{tree-like} & 13 & 24 & $\mathbf{-11}$ \\
\bottomrule
\end{tabular}
\end{center}

All three break earlier than a tree, \emph{including the negative control}. The
ordering is in the right direction --- the reticulate populations break earliest,
the tree-like one is intermediate, a genuine tree latest --- but the statistic
does not separate. Any threshold between $13$ and $24$ would classify all three
correctly, and selecting such a threshold from three populations with one
negative control is fitting a cutoff to the answer. We select none and report the
comparison as it stands. Both alignment references agree on every verdict, so
none of this is an alignment artefact.

\begin{remark}
One negative control cannot distinguish ``this method has a false-positive
problem'' from ``this population's tree-like prior is wrong.'' These priors are
priors: \cite{paperf} already records that one of the negative-control families
has documented within-type recombination~\cite{murahwa}, and a prior drawn from a
literature consensus is not a label. With a single control the two explanations
are observationally identical, and no additional statistic computed on the same
three populations can separate them.
\end{remark}

\section{Related work}\label{sec:related}

Split decomposition~\cite{bandeltdress} decomposes a metric into a totally
split-decomposable part and a split-prime residue, and the isolation index
of~\eqref{eq:iso} is its central quantity; the splits-equivalence theorem
relating pairwise compatibility to displayability by a tree is
standard~\cite{semplesteel}. Our contribution at this level is not the index but
the observation that complete deletion within the subset makes its defining
minimum incommensurable, and that a conjunction recovers an exact substitute.

Statistics for non-tree-likeness computed by aggregating over quadruples include
$\delta$ plots~\cite{deltaplots}, the pairwise homoplasy index~\cite{bruen}, and
quartet-based summary tests for introgression~\cite{frankelane}. These average or
count; the certificate exhibits. The distinction is the point of the programme
rather than a claim of superior power, and~\cite{paperf} reports the quartet
certificate tying the pairwise homoplasy index on power under matched nulls.

The relationship between split decomposability and the topology of
Vietoris--Rips complexes is actively studied: Gómez~\cite{gomez} investigates
Vietoris--Rips complexes of split-decomposable spaces, motivated by the fact that
totally decomposable spaces generalise metric trees and are therefore attractive
candidates for efficiently computable persistent homology. We note this because a
companion result closed a distance-magnitude Rips filtration on measured grounds
for this task, and that closure concerned a specific implementation --- a global
complex cannot use per-subset complete deletion, and a magnitude-driven
filtration tracks divergence rather than the index a tree forces to zero --- and
not the mathematical direction, which~\cite{gomez} shows to be live. We have read
only the abstract of~\cite{gomez} and cite it for that framing alone.

\section{Limitations}\label{sec:limits}

The formal content of \S\ref{sec:lift} is proved and the empirical content is
one simulated grid. Specifically:

\begin{itemize}
\item \textbf{The lift is unvalidated on real data.} It works at $8$--$10$
simulated taxa and does not transfer to $16$ real taxa. The strict form has no
power there and the calibrated form cannot be shown to separate.
\item \textbf{Three populations, one negative control.} This cannot establish
specificity, and we do not claim it does. The admissibility gate is what limits
the sample, not the availability of sequences.
\item \textbf{Theorem~\ref{thm:lift} is about three rationals.} It says nothing
about whether a certified split corresponds to a biological event; that inference
remains the reader's, informed by the exhibited sequences.
\item \textbf{The exception budget weakens the claim.} ``At most $k$ conjuncts
fail'' is exact but is a weaker statement than the strict conjunction, and its
calibration depends on the matched null being correctly matched --- which
\S\ref{sec:nullmatching} shows is a non-trivial requirement that an earlier
version of this apparatus did not meet.
\item \textbf{Priors are not labels.} The literature expectations used as ground
truth are consensus readings, one of which is known to be weak.
\end{itemize}

What would settle it is more admissible negative controls --- by a better
alignment step whose error is measured at the loads real data presents, or by
pre-aligned data of independent provenance --- and not another statistic computed
on these three populations.

\section{Reproduction}\label{sec:repro}

All numbers in this paper are produced by committed code and stored artifacts,
and no result is hand-entered. The Lean module is verified by compiling and
reading \texttt{\#print axioms} for every declaration. The simulated experiment,
the parameter sweeps, the nuisance-variability measurement and the real-data run
are separate subcommands writing JSON artifacts; the refactor that made the
$16$-taxon computation possible is confirmed behaviour-preserving against the
committed simulated artifact. The accompanying software release records the test
counts, the clippy warning count, and the guard outputs for the edition that
produced these figures.

\begin{thebibliography}{99}

\bibitem{buneman} Buneman P (1971/1974). The recovery of trees from measures of
dissimilarity. In \emph{Mathematics in the Archaeological and Historical
Sciences}, Edinburgh University Press, 387--395; and: A note on the metric
properties of trees. \emph{Journal of Combinatorial Theory} B 17, 48--50.

\bibitem{bandeltdress} Bandelt H-J, Dress AWM (1992). A canonical decomposition
theory for metrics on a finite set. \emph{Advances in Mathematics} 92(1),
47--105. \url{https://doi.org/10.1016/0001-8708(92)90061-O}

\bibitem{semplesteel} Semple C, Steel M (2003). \emph{Phylogenetics}. Oxford
University Press.

\bibitem{deltaplots} Holland BR, Huber KT, Dress A, Moulton V (2002). $\delta$
plots: a tool for analyzing phylogenetic distance data. \emph{Molecular Biology
and Evolution} 19(12), 2051--2059.

\bibitem{bruen} Bruen TC, Philippe H, Bryant D (2006). A simple and robust
statistical test for detecting the presence of recombination. \emph{Genetics}
172(4), 2665--2681.

\bibitem{frankelane} Frankel LE, An\'e C (2023). Summary tests of introgression
are highly sensitive to rate variation across lineages. \emph{Systematic
Biology} 72(5), 1057--1069.

\bibitem{murahwa} Murahwa AT, \emph{et al.}\ (2020). Recombination between
high-risk human papillomaviruses and non-human primate papillomaviruses:
evidence of ancient host switching. \emph{Virus Evolution} 6(2), veaa073.

\bibitem{gomez} G\'omez M (2024). Vietoris--Rips complexes of split-decomposable
spaces. arXiv:2403.15655. \url{https://arxiv.org/abs/2403.15655}

\bibitem{paperf} Callens X (2026). Proof-carrying detection of reticulate
evolution: a tree-level null, the measured cost of the alignment step, and the
Gromov hinge of Buneman's converse. Edition 7. Zenodo.
\url{https://doi.org/10.5281/zenodo.23602595}

\bibitem{textbook} Callens X (2026). \emph{Formal Genetics: Proof, Computation
and Verified Reasoning in Mathematical Genetics and Biophysics}. Edition 11.
Zenodo. \url{https://doi.org/10.5281/zenodo.23602616}

\end{thebibliography}

\end{document}
